\documentclass[12pt]{article}
\usepackage{geometry}
\geometry{margin=1in}
\usepackage{hyperref}

\begin{document}

\noindent
March 2026

\bigskip

\noindent
Professor Mariarosaria Taddeo\\
Editor-in-Chief\\
\textit{Minds and Machines}

\bigskip

\noindent
Dear Professor Taddeo,

\bigskip

I am writing to submit the manuscript entitled ``From Pathfinding to Field Theory: The Transition from $A^*$ on the Moral Manifold to the Moral Field Equation'' for consideration as an original research article in \textit{Minds and Machines}.

This paper articulates the mathematical transition from a computationally tractable model of moral reasoning ($A^*$ pathfinding on a moral manifold) to a full field theory, in which the curvature of moral space is determined by the distribution of moral commitments, relationships, and stakes---analogous to the transition from Newtonian mechanics to general relativity. The paper identifies four specific mathematical upgrades (dynamic metric tensor, moral Lagrangian, Riemann curvature tensor, and moral field equation), locates three of the four within existing proven results, and presents the fourth as the principal open problem.

A notable empirical contribution is the falsification of the originally proposed $SU(2)_I \times U(1)_H$ gauge group via CHSH Bell tests ($N = 600$), establishing the classical $D_4 \times U(1)_H$ as the correct symmetry structure. This demonstrates the framework's commitment to empirical falsifiability.

I believe this manuscript is well suited to \textit{Minds and Machines} for three reasons:

\begin{enumerate}
    \item \textbf{Philosophy-physics-computation intersection.} The paper bridges mathematical physics (gauge theory, Riemannian geometry, Lagrangian mechanics), philosophy (metaethics, deontic logic), and computation (AI alignment, $A^*$ search), which is the core territory of \textit{Minds and Machines}.

    \item \textbf{Honest epistemic status marking.} The paper distinguishes rigorously between what is proven (Noether conservation, gauge group structure), what is empirically confirmed (metric calibration, cross-lingual invariance), what has been falsified ($SU(2)$ contextuality), and what remains open (the moral field equation). This transparency is central to the paper's contribution.

    \item \textbf{Companion to a manuscript under review.} A foundational paper on the Geometric Ethics framework (``Geometric Ethics: Invariance, Conservation, and the Structure of Moral Evaluation'') is currently under review at \textit{Minds and Machines}. The present manuscript extends that work by developing the field-theoretic architecture that the foundational paper identifies as the principal open research direction. It is self-contained and does not depend on acceptance of the companion paper.
\end{enumerate}

This manuscript has not been published previously and is not under consideration elsewhere. The work is original, and I have read and comply with the journal's ethical guidelines.

I have suggested six reviewers with relevant expertise in the Editorial Manager system.

Thank you for considering this manuscript. I look forward to your response.

\bigskip

\noindent
Sincerely,

\bigskip

\noindent
Andrew H. Bond\\
Senior Lecturer\\
Department of Computer Engineering\\
San Jos\'{e} State University\\
\href{mailto:andrew.bond@sjsu.edu}{andrew.bond@sjsu.edu}

\end{document}
